%%%%%%%%%%%%%%%%
%%% deut 1 %%%
%%%%%%%%%%%%%%%%

\Chapter{1}

\Verse{1}
{
	\hDtr1{אֵלֶּה הַדְּבָרִים אֲשֶׁר דִּבֶּר מֹשֶׁה אֶל כּל יִשְׂרָאֵל בְּעֵבֶר הַיַּרְדֵּן בַּמִּדְבָּר בָּעֲרָבָה מוֹל סוּף בֵּין פָּארָן וּבֵין תֹּפֶל וְלָבָן וַחֲצֵרֹת וְדִי זָהָב׃}
}
{
	\eDtr1{
		These are the words that Moses spoke to all of Israel\\
		across the Jordan in the wilderness, in the plain opposite\\
		Suph between Paran and Tophel and Laban and Hazeroth and\\
		Di-zahab,\\
	}
}
{
%	\footnote{
%		words. The last verse in Numbers is, “These are the commandments and the
%		judgments that YHWH commanded by Moses’ hand to the children of Israel.”
%		The first verse in Deuteronomy is, “These are the words that Moses spoke
%		to all of Israel.” The difference between commandments and judgments in
%		Numbers and words in Deuteronomy marks an essential change. Commandments
%		and judgments are laws. Now, in the last book of the Torah, Moses does
%		more than give law. Deuteronomy is a book of Words. That is its name in
%		Hebrew: dèbãrîm, and that is what it is about. Its first thirty chapters
%		are Moses’ farewell address to his people before his death. It would
%		have taken close to three hours to say it all to them. It contains
%		history, law, and great wisdom. In places, especially near its end, it
%		is beautiful—inspired and inspirational. Moses is eloquent. And that is
%		ironic and instructive when we turn back to Moses’ first meeting with
%		God, at the burning bush. There he tries to escape from the assignment
%		to go speak to the Pharaoh by saying, “I’m not a man of words” (Exod
%		4:10)! Now he has become a man of words. It is interesting, remarkable,
%		ironic, and inspiring to see Moses’ development through all that has
%		happened in forty years into a man of words. More than any other human
%		in the Bible, Moses grows and changes in the course of his life. One can
%		change: change professions, change values, change lifestyle, change
%		character. One can grow and become stronger and better. This is also a
%		change in the presentation of law in the Torah. Most of the commandments
%		in the Torah have been given without reasons or explanations. From the
%		law of the “red cow” (Numbers 19) to the Ten Commandments, one is not
%		told why one must perform them, but only that God commands it (with a
%		few notable exceptions). But the law code of Deuteronomy (12–28) is
%		preceded by eleven chapters of history, explanation, and inspiration,
%		and it is followed by two chapters of exquisite revelation of the
%		relevance and value of the commandments. The Torah thus concludes with
%		the message that the law is meant to be relevant, comprehensible, and
%		meaningful in the people’s lives. It is appropriate to seek out the
%		meanings of the laws and, when interpreting the law, to understand that
%		it is explicitly meant to enhance lives. One must not apply it in a way
%		that causes injury or undermines its positive function in life. Footnote
%		1:1. Moses. In the Hebrew Bible’s picture of human history, Moses is the
%		first great man—and the first great leader. In Deuteronomy, our
%		acquaintance with Moses reaches a summit. Deuteronomy is the book that
%		is about Moses more than any other in the Torah. All are called “the
%		Five Books of Moses,” but Deuteronomy is the most focused on Moses
%		himself: a speech by him, expressions of his own views, descriptions of
%		his last acts. By virtue of the book’s being a speech, Moses now speaks
%		for himself, as it were, rather than the narrator’s speaking for and
%		about him. Moses has moved from “heavy mouth” to eloquence. Essentially
%		an unknown quantity at the burning bush, he has become probably the most
%		known person in the Bible. So known and still so enigmatic: the man who
%		wears a veil; who feared to take the assignment to go to Egypt but whose
%		temper flared at the king of Egypt; who ran from a staff that turned to
%		a snake at the burning bush but who disobediently struck a rock with a
%		staff at Meribah; who pleaded with God not to be angry with the people
%		over their golden calf but who smashed the tablets when he saw the calf
%		himself; the humblest man on earth who spoke to the deity more
%		audaciously than any other human. In analyzing and interpreting Moses,
%		one is more likely to reveal things about oneself than about Moses.
%		Moses is so much a Rorschach test because he is so much a regular human
%		being, with weaknesses, a temper, fears, and flaws—with key pieces
%		unrevealed—who comes to be the leader of a nation, a spokesman for God,
%		and the founder of a faith that plays a role in the destiny of
%		humankind. That is Moses as portrayed in the text. The reader’s
%		knowledge, moreover, that Moses is to some degree—to whatever
%		degree—historical behind this portrayal fuels this phenomenon. This
%		real, regular, unique, meek, powerful, audacious figure cries out for
%		understanding and interpretation. Footnote 1:1. the words that Moses
%		spoke. In developing from a man of actions to a man of words, Moses
%		imitates God. The Tanak depicts God as becoming more and more hidden
%		over the course of history. In the first books of the Bible God appears
%		to humans, is seen and heard at Sinai, makes His presence known through
%		miracles, angels, and the column of cloud and fire. But these visible
%		signs of divine action in history disappear from the story one by one.
%		And by the last books of the Tanak, there are no angels or miracles. The
%		words “YHWH appeared to” and “YHWH spoke to” do not occur to anyone.
%		Instead, the priest Ezra reads the Torah aloud to the people. In the
%		place of the acts of God there is the word of God. When the Torah
%		pictures Moses ending his life in words, he imitates and prefigures the
%		transformation of the human experience of God that will occur in the
%		Bible.
%	}
}

\Verse{2}
{
	\hDtr1{אַחַד עָשָׂר יוֹם מֵחֹרֵב דֶּרֶךְ הַר שֵׂעִיר עַד קָדֵשׁ בַּרְנֵעַ׃}
}
{
	\eDtr1{
		eleven days from Horeb by way of Mount Seir up to\\
		Kadesh-barnea.\\
	}
}
{
%	\footnote{
%		eleven days from Horeb. The journey from Horeb (the other name for Mount
%		Sinai) to the promised land only needed to take eleven days! Now, as the
%		people stand at the entrance to the land after forty years, after
%		everyone with adult memories of slavery has died, the meaning of the
%		loss of forty years can be deeply felt, appreciated, and regretted. The
%		new generation can think: “If only my mother and father had lived to see
%		this.” So they quintessentially reflect the human experience of wishing,
%		when we experience some joy after our parents or grandparents have died,
%		that they had lived to see it.
%	}
}

\Verse{3}
{
	\hDtr1{וַיְהִי בְּאַרְבָּעִים שָׁנָה בְּעַשְׁתֵּי עָשָׂר חֹדֶשׁ בְּאֶחָד לַחֹדֶשׁ דִּבֶּר מֹשֶׁה אֶל בְּנֵי יִשְׂרָאֵל כְּכֹל אֲשֶׁר צִוָּה יְהֹוָה אֹתוֹ אֲלֵהֶם׃}
}
{
	\eDtr1{
		And it was in the fortieth year, in the eleventh month, on\\
		the first of the month, Moses spoke to the children of\\
		Israel of everything that YHWH commanded him to them.\\
	}
}
{
%	\footnote{
%		in the fortieth year. The juxtaposition of this notice (that it was in
%		the fortieth year) to the preceding verse (that it was only eleven days’
%		journey) underscores the point that I made in my preceding comment.
%		These two verses should always be read together as a unit: “It is eleven
%		days from Horeb … but it was in the fortieth year …” And this
%		juxtaposition cries out for us to consider all the ways of interpreting
%		this fact: the irony, the waste, the necessity, the experience, the
%		growth, the suffering, the need for a new generation who had not known
%		slavery, the inability of humans to outgrow a trauma of mammoth
%		proportions to an entire nation (I write this in the generation after
%		the holocaust, in which the psychological effects of what happened still
%		afflict us).
%	}
}

\Verse{4}
{
	\hDtr1{אַחֲרֵי הַכֹּתוֹ אֵת סִיחֹן מֶלֶךְ הָאֱמֹרִי אֲשֶׁר יוֹשֵׁב בְּחֶשְׁבּוֹן וְאֵת עוֹג מֶלֶךְ הַבָּשָׁן אֲשֶׁר יוֹשֵׁב בְּעַשְׁתָּרֹת בְּאֶדְרֶעִי׃}
}
{
	\eDtr1{
		After he struck Sihon, king of the Amorites, who lived in\\
		Heshbon, and Og, king of Bashan, who lived in Ashtaroth at\\
		Edrei,\\
	}
}
{
}

\Verse{5}
{
	\hDtr1{בְּעֵבֶר הַיַּרְדֵּן בְּאֶרֶץ מוֹאָב הוֹאִיל מֹשֶׁה בֵּאֵר אֶת הַתּוֹרָה הַזֹּאת לֵאמֹר׃}
}
{
	\eDtr1{
		across the Jordan in the land of Moab, Moses undertook to\\
		make this instruction clear, saying,\\
	}
}
{
%	\footnote{
%		undertook. This is the word that Abraham uses twice in his attempt to
%		affect God’s decision about the fate of Sodom. He says: “Here I’ve
%		undertaken to speak to my Lord, and I’m dust and ashes” (Gen 18:27,31).
%		That is the first time in the Tanak that a human argues with God. Now
%		the same term is used for Moses’ setting out to convey a message from
%		God to the people. Again it is a reminder of how far humans have come,
%		and it may even hint that such growth was necessary: humans had to rise
%		to that level, at which one (Abraham) could question God, before they
%		could move to the level of receiving the commandments. Humans have to be
%		able to question commandments—even divine commandments, especially
%		divine commandments—in order for there to be any merit in performing
%		them. Footnote 1:5. make clear. Hebrew b[image]’[image]r. This term
%		occurs just twice in the Torah: here at the beginning of Moses’ speech
%		and, later, near the end of the speech (Deut 27:8). There it refers to
%		the way in which this instruction is to be written on stones: very
%		clearly. Here Moses seeks to accomplish in words what is accomplished
%		there in writing, to make the instructions clear. Among his final words
%		he will tell the people that this is not a hidden or distant thing. It
%		is close and is meant to be practiced (29:28; 30:11–14).
%	}
}

\Verse{6}
{
	\hDtr1{יְהֹוָה אֱלֹהֵינוּ דִּבֶּר אֵלֵינוּ בְּחֹרֵב לֵאמֹר רַב לָכֶם שֶׁבֶת בָּהָר הַזֶּה׃}
}
{
	\eDtr1{
		“YHWH our God spoke to us in Horeb, saying, ‘You’ve had\\
		enough of staying at this mountain.\\
	}
}
{
}

\Verse{7}
{
	\hDtr1{פְּנוּ וּסְעוּ לָכֶם וּבֹאוּ הַר הָאֱמֹרִי וְאֶל כּל שְׁכֵנָיו בָּעֲרָבָה בָהָר וּבַשְּׁפֵלָה וּבַנֶּגֶב וּבְחוֹף הַיָּם אֶרֶץ הַכְּנַעֲנִי וְהַלְּבָנוֹן עַד הַנָּהָר הַגָּדֹל נְהַר פְּרָת׃}
}
{
	\eDtr1{
		Turn and travel and come to the Amorite hill country and\\
		to all its neighboring places, in the plain, in the hill\\
		country, and in the lowland and in the Negeb and by the\\
		seashore, the land of the Canaanite and Lebanon as far as\\
		the big river, the Euphrates River.\\
	}
}
{
}

\Verse{8}
{
	\hDtr1{רְאֵה נָתַתִּי לִפְנֵיכֶם אֶת הָאָרֶץ בֹּאוּ וּרְשׁוּ אֶת הָאָרֶץ אֲשֶׁר נִשְׁבַּע יְהֹוָה לַאֲבֹתֵיכֶם לְאַבְרָהָם לְיִצְחָק וּלְיַעֲקֹב לָתֵת לָהֶם וּלְזַרְעָם אַחֲרֵיהֶם׃}
}
{
	\eDtr1{
		See: I’ve put the land in front of you. Come and possess\\
		the land that YHWH swore to your fathers, to Abraham, to\\
		Isaac, and to Jacob, to give to them and to their seed\\
		after them.’\\
	}
}
{
}

\Verse{9}
{
	\hDtr1{וָאֹמַר אֲלֵכֶם בָּעֵת הַהִוא לֵאמֹר לֹא אוּכַל לְבַדִּי שְׂאֵת אֶתְכֶם׃}
}
{
	\eDtr1{
		“And I said to you at that time, saying, ‘I’m not able to\\
		carry you by myself.\\
	}
}
{
%	\footnote{
%		I said to you at that time. Moses says, “I said to you at that time,”
%		even though most of the people he is addressing now were not even born
%		yet at the time he is discussing. And Moses will continue to speak this
%		way in this address. He conveys to the people that they are part of
%		history, that they receive the burden and the legacy of their parents
%		and ancestors. And in the course of his speech he will extend this
%		beyond his immediate audience as well. He mixes past, present, and
%		future generations. So in the end it becomes the message for each
%		generation who reads this text.
%	}
}

\Verse{10}
{
	\hDtr1{יְהֹוָה אֱלֹהֵיכֶם הִרְבָּה אֶתְכֶם וְהִנְּכֶם הַיּוֹם כְּכוֹכְבֵי הַשָּׁמַיִם לָרֹב׃}
}
{
	\eDtr1{
		YHWH, your God, has made you multiply, and here today\\
		you’re like the stars of the skies for multitude.\\
	}
}
{
}

\Verse{11}
{
	\hDtr1{יְהֹוָה אֱלֹהֵי אֲבוֹתֵכֶם יֹסֵף עֲלֵיכֶם כָּכֶם אֶלֶף פְּעָמִים וִיבָרֵךְ אֶתְכֶם כַּאֲשֶׁר דִּבֶּר לָכֶם׃}
}
{
	\eDtr1{
		(May YHWH, your fathers’ God, add on to you a thousand\\
		times more like you! And may He bless you as He spoke to\\
		you.)\\
	}
}
{
}

\Verse{12}
{
	\hDtr1{אֵיכָה אֶשָּׂא לְבַדִּי טרְחֲכֶם וּמַשַּׂאֲכֶם וְרִיבְכֶם׃}
}
{
	\eDtr1{
		How shall I carry your stress and your burden and your\\
		quarrels by myself?\\
	}
}
{
}

\Verse{13}
{
	\hDtr1{הָבוּ לָכֶם אֲנָשִׁים חֲכָמִים וּנְבֹנִים וִידֻעִים לְשִׁבְטֵיכֶם וַאֲשִׂימֵם בְּרָאשֵׁיכֶם׃}
}
{
	\eDtr1{
		Get wise and understanding and knowledgeable people for\\
		your tribes, and I’ll set them among your heads.’\\
	}
}
{
}

\Verse{14}
{
	\hDtr1{וַתַּעֲנוּ אֹתִי וַתֹּאמְרוּ טוֹב הַדָּבָר אֲשֶׁר דִּבַּרְתָּ לַעֲשׂוֹת׃}
}
{
	\eDtr1{
		“And you answered me, and you said, ‘The thing that you’ve\\
		spoken is good to do.’\\
	}
}
{
}

\Verse{15}
{
	\hDtr1{וָאֶקַּח אֶת רָאשֵׁי שִׁבְטֵיכֶם אֲנָשִׁים חֲכָמִים וִידֻעִים וָאֶתֵּן אוֹתָם רָאשִׁים עֲלֵיכֶם שָׂרֵי אֲלָפִים וְשָׂרֵי מֵאוֹת וְשָׂרֵי חֲמִשִּׁים וְשָׂרֵי עֲשָׂרֹת וְשֹׁטְרִים לְשִׁבְטֵיכֶם׃}
}
{
	\eDtr1{
		“And I took heads of your tribes, wise and knowledgeable\\
		people, and I set them as heads over you, chiefs of\\
		thousands and chiefs of hundreds and chiefs of fifties and\\
		chiefs of tens and officers for your tribes.\\
	}
}
{
}

\Verse{16}
{
	\hDtr1{וָאֲצַוֶּה אֶת שֹׁפְטֵיכֶם בָּעֵת הַהִוא לֵאמֹר שָׁמֹעַ בֵּין אֲחֵיכֶם וּשְׁפַטְתֶּם צֶדֶק בֵּין אִישׁ וּבֵין אָחִיו וּבֵין גֵּרוֹ׃}
}
{
	\eDtr1{
		And I commanded your judges at that time, saying, ‘Hear\\
		between your brothers, and you shall judge with justice\\
		between a man and his brother or between him and his\\
		alien.\\
	}
}
{
%	\footnote{
%		Hear between your brothers. “Hear” is meant in the same sense as the
%		phrase in English: judges “hear a case.” Footnote 1:16. his alien. The
%		possessive (his alien) is used here and in the Ten Commandments (your
%		alien; Exod 20:10; Deut 5:14), conveying how one must feel about—and
%		treat—a non-Israelite living among the people of Israel.
%	}
}

\Verse{17}
{
	\hDtr1{לֹא תַכִּירוּ פָנִים בַּמִּשְׁפָּט כַּקָּטֹן כַּגָּדֹל תִּשְׁמָעוּן לֹא תָגוּרוּ מִפְּנֵי אִישׁ כִּי הַמִּשְׁפָּט לֵאלֹהִים הוּא וְהַדָּבָר אֲשֶׁר יִקְשֶׁה מִכֶּם תַּקְרִבוּן אֵלַי וּשְׁמַעְתִּיו׃}
}
{
	\eDtr1{
		You shall not recognize a face in judgment: you shall hear\\
		the same for small and for big. Don’t be fearful in front\\
		of a man—because justice is God’s. And a matter that will\\
		be too hard for you, you shall bring forward to me, and\\
		I’ll hear it.’\\
	}
}
{
%	\footnote{
%		recognize a face in judgment. Meaning: you shall not show favoritism to
%		anyone in a legal proceeding.
%	}
}

\Verse{18}
{
	\hDtr1{וָאֲצַוֶּה אֶתְכֶם בָּעֵת הַהִוא אֵת כּל הַדְּבָרִים אֲשֶׁר תַּעֲשׂוּן׃}
}
{
	\eDtr1{
		And I commanded you at that time all the things that you\\
		should do.\\
	}
}
{
}

\Verse{19}
{
	\hDtr1{וַנִּסַּע מֵחֹרֵב וַנֵּלֶךְ אֵת כּל הַמִּדְבָּר הַגָּדוֹל וְהַנּוֹרָא הַהוּא אֲשֶׁר רְאִיתֶם דֶּרֶךְ הַר הָאֱמֹרִי כַּאֲשֶׁר צִוָּה יְהֹוָה אֱלֹהֵינוּ אֹתָנוּ וַנָּבֹא עַד קָדֵשׁ בַּרְנֵעַ׃}
}
{
	\eDtr1{
		“And we traveled from Horeb and went through all that big\\
		and fearful wilderness that you’ve seen by way of the\\
		Amorite hill country as YHWH, our God, commanded us, and\\
		we came to Kadesh-barnea.\\
	}
}
{
%	\footnote{
%		from Horeb … to Kadesh-barnea. Moses covers the trip from Horeb to the
%		area near the promised land in a single sentence. Thus the text again
%		conveys that the journey could have been a simple thing—eleven days, as
%		in v. 2—but the people’s condition turned it into forty years.
%	}
}

\Verse{20}
{
	\hDtr1{וָאֹמַר אֲלֵכֶם בָּאתֶם עַד הַר הָאֱמֹרִי אֲשֶׁר יְהֹוָה אֱלֹהֵינוּ נֹתֵן לָנוּ׃}
}
{
	\eDtr1{
		And I said to you, ‘You’ve come to the Amorite hill\\
		country that YHWH, our God, is giving us.\\
	}
}
{
}

\Verse{21}
{
	\hDtr1{רְאֵה נָתַן יְהֹוָה אֱלֹהֶיךָ לְפָנֶיךָ אֶת הָאָרֶץ עֲלֵה רֵשׁ כַּאֲשֶׁר דִּבֶּר יְהֹוָה אֱלֹהֵי אֲבֹתֶיךָ לָךְ אַל תִּירָא וְאַל תֵּחָת׃}
}
{
	\eDtr1{
		See: YHWH, our God, has put the land in front of you. Go\\
		up, possess it as YHWH, your fathers’ God, spoke to you.\\
		Don’t be afraid and don’t be dismayed.’\\
	}
}
{
%	\footnote{
%		See.... Go up, possess. Not only does Moses mix past, present, and
%		future generations in his wording. Now he mixes the whole community and
%		each individual. Until now he has spoken in the plural, to all of Israel
%		in front of him. But now he changes to the singular, telling each
%		person, young and old, male and female: YHWH has put this land in front
%		of you.
%	}
}

\Verse{22}
{
	\hDtr1{וַתִּקְרְבוּן אֵלַי כֻּלְּכֶם וַתֹּאמְרוּ נִשְׁלְחָה אֲנָשִׁים לְפָנֵינוּ וְיַחְפְּרוּ לָנוּ אֶת הָאָרֶץ וְיָשִׁבוּ אֹתָנוּ דָּבָר אֶת הַדֶּרֶךְ אֲשֶׁר נַעֲלֶה בָּהּ וְאֵת הֶעָרִים אֲשֶׁר נָבֹא אֲלֵיהֶן׃}
}
{
	\eDtr1{
		“And you came forward to me, all of you, and said, ‘Let’s\\
		send men ahead of us, and let them explore the land for us\\
		and bring us back word of the way by which we’ll go up and\\
		of the cities to which we’ll come.’\\
	}
}
{
}

\Verse{23}
{
	\hDtr1{וַיִּיטַב בְּעֵינַי הַדָּבָר וָאֶקַּח מִכֶּם שְׁנֵים עָשָׂר אֲנָשִׁים אִישׁ אֶחָד לַשָּׁבֶט׃}
}
{
	\eDtr1{
		And the thing was good in my eyes, and I took twelve men\\
		from you, one man per tribe.\\
	}
}
{
}

\Verse{24}
{
	\hDtr1{וַיִּפְנוּ וַיַּעֲלוּ הָהָרָה וַיָּבֹאוּ עַד נַחַל אֶשְׁכֹּל וַיְרַגְּלוּ אֹתָהּ׃}
}
{
	\eDtr1{
		And they turned and went up to the hill country and came\\
		to the Wadi Eshcol and spied it out.\\
	}
}
{
}

\Verse{25}
{
	\hDtr1{וַיִּקְחוּ בְיָדָם מִפְּרִי הָאָרֶץ וַיּוֹרִדוּ אֵלֵינוּ וַיָּשִׁבוּ אֹתָנוּ דָבָר וַיֹּאמְרוּ טוֹבָה הָאָרֶץ אֲשֶׁר יְהֹוָה אֱלֹהֵינוּ נֹתֵן לָנוּ׃}
}
{
	\eDtr1{
		And they took some of the land’s fruit in their hand and\\
		brought it down to us, and they brought us back word and\\
		said, the land that YHWH, our God, is giving us is good.\\
	}
}
{
}

\Verse{26}
{
	\hDtr1{וְלֹא אֲבִיתֶם לַעֲלֹת וַתַּמְרוּ אֶת פִּי יְהֹוָה אֱלֹהֵיכֶם׃}
}
{
	\eDtr1{
		But you weren’t willing to go up, and you rebelled against\\
		the word of YHWH, your God,\\
	}
}
{
}

\Verse{27}
{
	\hDtr1{וַתֵּרָגְנוּ בְאהֳלֵיכֶם וַתֹּאמְרוּ בְּשִׂנְאַת יְהֹוָה אֹתָנוּ הוֹצִיאָנוּ מֵאֶרֶץ מִצְרָיִם לָתֵת אֹתָנוּ בְּיַד הָאֱמֹרִי לְהַשְׁמִידֵנוּ׃}
}
{
	\eDtr1{
		and you grumbled in your tents and said, ‘Because of\\
		YHWH’s hatred for us He brought us out from the land of\\
		Egypt, to put us in the Amorite’s hand to destroy us!\\
	}
}
{
}

\Verse{28}
{
	\hDtr1{אָנָה אֲנַחְנוּ עֹלִים אַחֵינוּ הֵמַסּוּ אֶת לְבָבֵנוּ לֵאמֹר עַם גָּדוֹל וָרָם מִמֶּנּוּ עָרִים גְּדֹלֹת וּבְצוּרֹת בַּשָּׁמָיִם וְגַם בְּנֵי עֲנָקִים רָאִינוּ שָׁם׃}
}
{
	\eDtr1{
		To where are we going up? Our brothers have melted our\\
		heart, saying, “We saw a bigger and taller people than we\\
		are, big cities and fortified to the skies, and also\\
		giants!”’\\
	}
}
{
}

\Verse{29}
{
	\hDtr1{וָאֹמַר אֲלֵכֶם לֹא תַעַרְצוּן וְלֹא תִירְאוּן מֵהֶם׃}
}
{
	\eDtr1{
		“And I said to you, ‘Don’t be scared and don’t be afraid\\
		of them.\\
	}
}
{
}

\Verse{30}
{
	\hDtr1{יְהֹוָה אֱלֹהֵיכֶם הַהֹלֵךְ לִפְנֵיכֶם הוּא יִלָּחֵם לָכֶם כְּכֹל אֲשֶׁר עָשָׂה אִתְּכֶם בְּמִצְרַיִם לְעֵינֵיכֶם׃}
}
{
	\eDtr1{
		YHWH, your God, who is going in front of you, He will\\
		fight for you—like everything that He did with you in\\
		Egypt before your eyes\\
	}
}
{
}

\Verse{31}
{
	\hDtr1{וּבַמִּדְבָּר אֲשֶׁר רָאִיתָ אֲשֶׁר נְשָׂאֲךָ יְהֹוָה אֱלֹהֶיךָ כַּאֲשֶׁר יִשָּׂא אִישׁ אֶת בְּנוֹ בְּכל הַדֶּרֶךְ אֲשֶׁר הֲלַכְתֶּם עַד בֹּאֲכֶם עַד הַמָּקוֹם הַזֶּה׃}
}
{
	\eDtr1{
		and in the wilderness, as you’ve seen that YHWH, your God,\\
		carried you the way a man would carry his child through\\
		all the way that you’ve gone until you came to this place.\\
	}
}
{
}

\Verse{32}
{
	\hDtr1{וּבַדָּבָר הַזֶּה אֵינְכֶם מַאֲמִינִם בַּיהֹוָה אֱלֹהֵיכֶם׃}
}
{
	\eDtr1{And in this thing, don’t you trust in YHWH, your God,}
}
{
}

\Verse{33}
{
	\hDtr1{הַהֹלֵךְ לִפְנֵיכֶם בַּדֶּרֶךְ לָתוּר לָכֶם מָקוֹם לַחֲנֹתְכֶם בָּאֵשׁ לַיְלָה לַרְאֹתְכֶם בַּדֶּרֶךְ אֲשֶׁר תֵּלְכוּ בָהּ וּבֶעָנָן יוֹמָם׃}
}
{
	\eDtr1{
		who was going in front of you on the way to scout a\\
		camping place for you, in fire at night to let you see the\\
		way in which you would go, and in a cloud by day?!’\\
	}
}
{
}

\Verse{34}
{
	\hDtr1{וַיִּשְׁמַע יְהֹוָה אֶת קוֹל דִּבְרֵיכֶם וַיִּקְצֹף וַיִּשָּׁבַע לֵאמֹר׃}
}
{
	\eDtr1{
		“And YHWH heard the sound of your words and was angry and\\
		swore, saying,\\
	}
}
{
}

\Verse{35}
{
	\hDtr1{אִם יִרְאֶה אִישׁ בָּאֲנָשִׁים הָאֵלֶּה הַדּוֹר הָרָע הַזֶּה אֵת הָאָרֶץ הַטּוֹבָה אֲשֶׁר נִשְׁבַּעְתִּי לָתֵת לַאֲבֹתֵיכֶם׃}
}
{
	\eDtr1{
		‘Not a man of these people, this bad generation, will see\\
		the good land that I swore to give to your fathers;\\
	}
}
{
}

\Verse{36}
{
	\hDtr1{זוּלָתִי כָּלֵב בֶּן יְפֻנֶּה הוּא יִרְאֶנָּה וְלוֹ אֶתֵּן אֶת הָאָרֶץ אֲשֶׁר דָּרַךְ בָּהּ וּלְבָנָיו יַעַן אֲשֶׁר מִלֵּא אַחֲרֵי יְהֹוָה׃}
}
{
	\eDtr1{
		just Caleb son of Jephunneh: he will see it, and I shall\\
		give the land on which he went to him and to his children\\
		because he went after YHWH fully.’\\
	}
}
{
}

\Verse{37}
{
	\hDtr1{גַּם בִּי הִתְאַנַּף יְהֹוָה בִּגְלַלְכֶם לֵאמֹר גַּם אַתָּה לֹא תָבֹא שָׁם׃}
}
{
	\eDtr1{
		(YHWH was incensed at me, too, because of you, saying,\\
		‘You, too, shall not come there.\\
	}
}
{
%	\footnote{
%		me, too. This is the first of many times that Moses will mention the
%		fact that he, too, will not live to see the land. Here he connects it
%		with the people’s actions, even though the reason that he cannot enter
%		the land is because of his sin at the rock in Meribah (Numbers 20). Even
%		though his comment comes in the context of the spies episode, he may
%		still be understood to be referring to the Meribah episode in this
%		comment. His point is to let the Israelites know that he shares their
%		fate. Like everyone else of his generation, he will not live to set foot
%		in the promised land. Footnote 1:37. because of you. If he is referring
%		to the Meribah episode, then he is blaming his actions there on the
%		people, whose rebellion moved him to act as he did. Alternatively, if he
%		is referring to the spies episode, then it is particularly striking that
%		Moses says “because of you,” since none of the people standing in front
%		of him was culpable for that sin. They were all the “infants” at the
%		time (v. 39), who are now adults; and they are to inherit the land
%		precisely because they were not part of the generation that rejected the
%		land in the spies episode. The point may be that Moses is especially
%		emphasizing the message that they are still a part of the history that
%		preceded them. Or this may convey that the concept of a “nation” or
%		“people” is a fluid one, crossing lines of generations. Like a river,
%		which is constantly being made up of new molecules of water yet is
%		always the same river (“You cannot step into the same river twice”), so
%		a people retains its identity even though its individual members are
%		constantly changing over time.
%	}
}

\Verse{38}
{
	\hDtr1{יְהוֹשֻׁעַ בִּן נוּן הָעֹמֵד לְפָנֶיךָ הוּא יָבֹא שָׁמָּה אֹתוֹ חַזֵּק כִּי הוּא יַנְחִלֶנָּה אֶת יִשְׂרָאֵל׃}
}
{
	\eDtr1{
		Joshua son of Nun, who is standing in front of you: he\\
		shall come there. Strengthen him, because he shall get\\
		Israel its legacy.’)\\
	}
}
{
}

\Verse{39}
{
	\hDtr1{וְטַפְּכֶם אֲשֶׁר אֲמַרְתֶּם לָבַז יִהְיֶה וּבְנֵיכֶם אֲשֶׁר לֹא יָדְעוּ הַיּוֹם טוֹב וָרָע הֵמָּה יָבֹאוּ שָׁמָּה וְלָהֶם אֶתְּנֶנָּה וְהֵם יִירָשׁוּהָ׃}
}
{
	\eDtr1{
		‘And your infants whom you said would become a spoil, and\\
		your children today who haven’t known good and bad, they\\
		will come there, and I’ll give it to them, and they will\\
		possess it.\\
	}
}
{
%	\footnote{
%		haven’t known good and bad. Elsewhere in the Tanak this expression
%		refers to a small child who is not yet old enough to know the difference
%		between good and bad (Isa 7:14–16). The reference to the knowledge of
%		good and bad is also one of many reminiscences here in the Torah’s last
%		book of things that occur at the beginning of the Torah’s first book. In
%		Genesis humans acquire knowledge of good and bad. Now Moses quotes God
%		from the spies episode, giving a clearer explanation of why only the
%		youngsters will live to enter the promised land: because they have not
%		yet “known good and bad.” Again knowledge of good and bad is tied up
%		with life and death.
%	}
}

\Verse{40}
{
	\hDtr1{וְאַתֶּם פְּנוּ לָכֶם וּסְעוּ הַמִּדְבָּרָה דֶּרֶךְ יַם סוּף׃}
}
{
	\eDtr1{
		But you: turn and travel to the wilderness by way of the\\
		Red Sea.’ 1:40; by way of the Red Sea. The Israelites\\
		never go back near any place that could be the site of the\\
		splitting of the sea. This reference therefore must be to\\
		the eastern arm of the Red Sea. It is further confirmation\\
		that the name of that entire body of water, Hebrew yam\\
		sûp, in the Torah refers to what is known today as the Red\\
		Sea, and not to any imagined “Reed Sea.” See the comments\\
		on Exod 13:8 and Num 21:4.\\
	}
}
{
}

\Verse{41}
{
	\hDtr1{וַתַּעֲנוּ וַתֹּאמְרוּ אֵלַי חָטָאנוּ לַיהֹוָה אֲנַחְנוּ נַעֲלֶה וְנִלְחַמְנוּ כְּכֹל אֲשֶׁר צִוָּנוּ יְהֹוָה אֱלֹהֵינוּ וַתַּחְגְּרוּ אִישׁ אֶת כְּלֵי מִלְחַמְתּוֹ וַתָּהִינוּ לַעֲלֹת הָהָרָה׃}
}
{
	\eDtr1{
		“And you answered and said to me, ‘We’ve sinned against\\
		YHWH. We’ll go up! And we’ll fight, according to\\
		everything that YHWH, our God, commanded us.’ And you each\\
		put on his weapons of war, and you were willing to go up\\
		to the hill country.\\
	}
}
{
%	\footnote{
%		you were willing. This verb (Hebrew t[image]hînû) occurs only here. No
%		one is sure what it means. Rashi connects it with the text of the spies
%		episode, in which the people say, “Here we are (Hebrew hinnennû), and
%		we’ll go up” (Num 14:40), which appears to be what Moses alludes to
%		here: “you were willing to go up.” That seems most probable to me, in
%		which case the text here makes a verb out of the word hinnennû (here we
%		are), conveying one’s willingness to do something.
%	}
}

\Verse{42}
{
	\hDtr1{וַיֹּאמֶר יְהֹוָה אֵלַי אֱמֹר לָהֶם לֹא תַעֲלוּ וְלֹא תִלָּחֲמוּ כִּי אֵינֶנִּי בְּקִרְבְּכֶם וְלֹא תִּנָּגְפוּ לִפְנֵי אֹיְבֵיכֶם׃}
}
{
	\eDtr1{
		“And YHWH said to me, ‘Say to them: You shall not go up,\\
		and you shall not fight, because I am not among you, so\\
		you won’t be stricken before your enemies.’\\
	}
}
{
}

\Verse{43}
{
	\hDtr1{וָאֲדַבֵּר אֲלֵיכֶם וְלֹא שְׁמַעְתֶּם וַתַּמְרוּ אֶת פִּי יְהֹוָה וַתָּזִדוּ וַתַּעֲלוּ הָהָרָה׃}
}
{
	\eDtr1{
		And I spoke to you, but you didn’t listen, and you\\
		rebelled against YHWH’s word and acted presumptuously and\\
		went up to the hill country. 1:43, didn’t listen. Moses\\
		reminds the people that “You didn’t listen,” and so later\\
		“YHWH didn’t listen.” This foreshadows developments coming\\
		in the Tanak in which God will be described as hiding His\\
		face, not seeing, and not listening—always in response to\\
		the people’s not listening to God. It will be stated\\
		explicitly among God’s last words to Moses at the end of\\
		the Torah (Deut 31:16–18).\\
	}
}
{
}

\Verse{44}
{
	\hDtr1{וַיֵּצֵא הָאֱמֹרִי הַיֹּשֵׁב בָּהָר הַהוּא לִקְרַאתְכֶם וַיִּרְדְּפוּ אֶתְכֶם כַּאֲשֶׁר תַּעֲשֶׂינָה הַדְּבֹרִים וַיַּכְּתוּ אֶתְכֶם בְּשֵׂעִיר עַד חרְמָה׃}
}
{
	\eDtr1{
		And the Amorite who lives in that hill country came out at\\
		you and pursued you the way bees do, and they crushed you\\
		in Seir as far as Hormah.\\
	}
}
{
}

\Verse{45}
{
	\hDtr1{וַתָּשֻׁבוּ וַתִּבְכּוּ לִפְנֵי יְהֹוָה וְלֹא שָׁמַע יְהֹוָה בְּקֹלְכֶם וְלֹא הֶאֱזִין אֲלֵיכֶם׃}
}
{
	\eDtr1{
		And you came back and wept in front of YHWH, but YHWH\\
		didn’t listen to your voices and didn’t hear you! 1:45,\\
		didn’t listen. Moses reminds the people that “You didn’t\\
		listen,” and so later “YHWH didn’t listen.” This\\
		foreshadows developments coming in the Tanak in which God\\
		will be described as hiding His face, not seeing, and not\\
		listening—always in response to the people’s not listening\\
		to God. It will be stated explicitly among God’s last\\
		words to Moses at the end of the Torah (Deut 31:16–18).\\
	}
}
{
}

\Verse{46}
{
	\hDtr1{וַתֵּשְׁבוּ בְקָדֵשׁ יָמִים רַבִּים כַּיָּמִים אֲשֶׁר יְשַׁבְתֶּם׃}
}
{
	\eDtr1{
		“So you lived in Kadesh many days—as many days as you\\
		lived there—\\
	}
}
{
}
